\section{相关工作}
\label{sec:related-work}

\subsection{自由手三维超声重建}

% P1：与英文 P1 对应。
自由手三维超声重建将不同探头位姿下采集的二维 B-mode 图像组合成一个体数据~\cite{prager2010three}。在带跟踪装置的系统中，探头定位与图像到探头的空间标定将各像素映射到统一的空间坐标系。随后，重建过程既需要对采样平面之间的区域进行插值，也需要融合相互重叠的观测。这些操作受到采集轨迹、空间采样和标定精度的影响~\cite{solberg2007freehand}。Rohling 等人~\cite{rohling1999interpolation}研究了径向基函数近似，并将其与三种标准重建方法进行比较，展示了在 B 型扫查切片间距不规则时重建质量与计算成本之间的权衡。

% P2：与英文 P2 对应。
连续插值函数可以给出已采集帧之间的数值，但边界保真度还取决于待融合观测之间的关系。配准误差可能使同一界面的观测落在不同空间位置，而声束入射方向的变化可能改变该界面的图像外观。这两类因素分别推动了几何校正和超声图像形成模型的研究~\cite{dou2026gau,wysocki2024ultranerf}。对于单一 B-mode 灰度场，观测边界的保持因此既取决于空间对齐，也取决于不同帧之间图像外观的一致性。

\subsection{超声隐式表示}

% P3：与英文 P3 对应。
隐式神经表示将信号参数化为坐标的函数。NeRF 将这种表示与渲染模型结合，用于新视角图像合成~\cite{mildenhall2020nerf}；M\"uller 等人~\cite{mueller2022instant}提出的多分辨率哈希编码则在多个空间尺度上提供可学习特征，降低神经图形基元的拟合成本。超声研究也已将坐标式表示用于体重建，例如 Song 等人~\cite{song2022carotid}针对颈动脉血管的研究。坐标式体模型支持在不同输出网格上重新采样~\cite{yeung2024sensorless}，但能够恢复的空间细节仍受到采集数据和重建模型的约束。

% P4：与英文 P4 对应。
神经表示在重建中承担的作用，还取决于采集几何是否可得及其可靠性。Yeung 等人~\cite{yeung2024sensorless}针对无外部跟踪传感器的胎脑重建，先估计图像帧的空间位置，再将这些位置与隐式体表示共同优化。针对带跟踪装置的自由手超声，ImplicitCell 预印本将超声分辨单元模型与体重建、位姿修正联合起来~\cite{song2025implicitcell}。GAU-NeRF 则通过联合优化场表示和探头位姿处理多次扫查之间的错配，并在优化早期利用梯度重加权和平滑提高优化稳定性~\cite{dou2026gau}。这些方法将图像帧的几何关系视为变量，与重建体联合估计。

% P5：与英文 P5 对应。
另一个问题是如何描述超声观测的形成过程。Ultra-NeRF 预测随空间变化的组织参数，并通过基于声线路径的超声渲染器描述衰减、反射和散射~\cite{wysocki2024ultranerf}。尽管其参数场本身按空间位置查询，图像形成过程中的方向依赖仍可通过渲染实现。UlRe-NeRF 预印本进一步研究了反射方向参数化和方向编码~\cite{guo2024ulre}。这类方法分别建模介质表示和观测图像合成，使探头方向能够影响渲染得到的 B-mode 图像外观。

% P6：与英文 P6 对应。
多尺度特征表示也已用于具有物理依据的超声渲染。Liu 等人~\cite{liu2026geometry}估计声阻抗，并利用其空间梯度对与几何相关的反射进行建模。该方法分别用粗尺度哈希特征表示反射、细尺度特征表示散射。UltraSoundNeRF 通过修订前向模型减少声学参数化中的冗余，并利用物理启发的正则化改善恢复所得声学参数图的可解释性~\cite{wysocki2026usnerf}。这些方法通过显式描述介质及其与超声的相互作用，对 B-mode 图像细节进行建模。

\subsection{结构先验与边界保持}

% P7：与英文 P7 对应。
超声图像滤波长期关注如何在抑制散斑的同时保持边界。散斑抑制各向异性扩散根据局部图像统计量调整扩散过程，使平滑操作能够响应散斑图像中的边缘~\cite{yu2002srad}。定向散斑抑制各向异性扩散进一步引入方向相关的滤波，在跨越轮廓的方向和沿主曲率方向采用不同的扩散强度~\cite{krissian2007osrad}。这些方法说明了局部图像结构如何指导随空间位置变化的正则化。它们主要面向图像或体数据的滤波；将类似的结构信息用于多帧连续场拟合时，还需要考虑不同观测之间的一致性。

% P8：与英文 P8 对应。
结构张量正则化在图像邻域内汇集梯度信息。Lefkimmiatis 等人~\cite{lefkimmiatis2015stv}提出结构张量全变分（STV），通过张量的特征值定义用于成像逆问题的一组正则项。Lefkimmiatis 和 Osher~\cite{lefkimmiatis2015nonlocal}进一步将其扩展为非局部结构张量泛函，纳入相似图像邻域的信息。由此得到的正则项结合局部方向结构与图像的非局部相似性，为从观测灰度中提取结构引导信息提供了依据。

% P9：与英文 P9 对应。
在超声领域，Li 等人~\cite{li2026ipbnlstv}已将非局部结构张量全变分（NLSTV）引入波束形成的逆问题表述。该方法对从射频测量数据重建复数超声图像的过程施加正则化，并在平面波成像数据上进行了评价。在这一任务中，先验作用于由声学测量数据形成图像的过程。B-mode 体重建则融合不同探头位姿下已经形成的图像。

% P10：与英文 P10 对应，保留研究定位的限定范围。
NeUF 在固定采集几何下，利用结构引导拟合连续的 B-mode 体表示。由训练帧计算得到的固定 NLSTV 结构响应用于监督辅助响应场。预测响应在最终灰度激活之前，对 logit 空间中的有界残差进行门控；观测 B-mode 灰度则作为重建目标。这种建模方式利用图像域先验调节直接拟合的灰度场中的细节，研究重点是观测图像边界的保真度。
